\originalchapter{课程作业: Lorentz 几何与共形能量}
\originalsection{作业7: 2011年11月3日交}
\begin{exercise}\label{ps:7I}
\textbf{原题 I.} 三维波方程的光滑初值 $f,g$ 支撑包含于闭球 $\overline B_R$, 证明 $|u(t,x)|\le C/(1+t)$, $C$ 仅依赖 $R,\|f\|_\infty,\|\nabla f\|_\infty,\|g\|_\infty$.
\end{exercise}
\begin{solution}
Kirchhoff 式写成
\[
u=\frac1{4\pi}\int_{\Sph^2}[f(x+t\omega)+t\nabla f(x+t\omega)\cdot\omega+tg(x+t\omega)]\dd\omega.
\]
$t\le2R$ 时取全球面积给 $|u|\le\|f\|_\infty+2R(\|\nabla f\|_\infty+\|g\|_\infty)$, 将常数乘 $1+2R$ 即满足所需界.

$t\ge2R$ 且移动球面与 $B_R$ 相交, 令 $d=|x|$, 则 $t-R\le d\le t+R$, 故 $d\ge t/2$. 球冠面积为 $2\pi t^2(1-\cos\theta_0)$, 其中 $\cos\theta_0=(d^2+t^2-R^2)/(2dt)$. 所以面积 $=\pi t[R^2-(d-t)^2]/d\le2\pi R^2$. 不交时面积0. 由 $\dd\omega=t^{-2}\dd S$, 得 $|u|\le (R^2/2t^2)\|f\|_\infty+(R^2/2t)(\|\nabla f\|_\infty+\|g\|_\infty)$. 在 $t\ge2R$ 把它改写为 $C/(1+t)$ 即可, $R>0$ 固定.
\end{solution}
\begin{exercise}\label{ps:7II}
\textbf{原题 II(a--c).} $\Lambda^Tm\Lambda=m$, $\Omega^Tm\Omega=m$, $m=\operatorname{diag}(-1,1,\ldots,1)$. 证明 $|\det\Lambda|=1$, 乘积与逆仍为 Lorentz 变换.
\end{exercise}
\begin{solution}
(a) 取行列式得 $(\det\Lambda)^2\det m=\det m$, 因 $\det m=-1\ne0$, 结论成立并保证可逆.
(b) $(\Lambda\Omega)^Tm(\Lambda\Omega)=\Omega^T(\Lambda^Tm\Lambda)\Omega=m$.
(c) 原等式两侧左乘 $(\Lambda^{-1})^T$、右乘 $\Lambda^{-1}$ 得 $m=(\Lambda^{-1})^Tm\Lambda^{-1}$.
\end{solution}
\begin{exercise}\label{ps:7III}
\textbf{原题 III.} $|v|<1$, $\gamma=(1-v^2)^{-1/2}$, 证明矩阵 $B$ 的 $(t,x^1)$ 块为 $\left(\begin{smallmatrix}\gamma&-\gamma v\\-\gamma v&\gamma\end{smallmatrix}\right)$、其它空间方向为恒等的 boost 为 Lorentz 变换.
\end{exercise}
\begin{solution}
二阶块计算给 $-(\gamma t-\gamma vx)^2+(\gamma x-\gamma vt)^2=\gamma^2(1-v^2)(-t^2+x^2)=-t^2+x^2$. 其它平方不变. 二次型等式极化后给 $B^TmB=m$. 二阶块行列式 $\gamma^2(1-v^2)=1$, 故总行列式1.
\end{solution}
\begin{exercise}\label{ps:7IV}
\textbf{原题 IV.} 未来类时向量 $X$ 可由行列式1的 Lorentz 变换变为 $(c,0,\ldots,0)$, $c>0$.
\end{exercise}
\begin{solution}
$n\ge2$ 时空间正旋转把 $\vec X$ 变为 $(a,0,\ldots)$, $a=|\vec X|$. $n=1$ 无须旋转, 用带符号的 $a=X^1$. 置 $v=a/X^0$, 类时性给 $|v|<1$. 上题 boost 后空间首分量 $\gamma(a-vX^0)=0$, 时间分量 $\gamma(X^0-va)=\sqrt{(X^0)^2-a^2}>0$. 旋转与 boost 均行列式1, 乘积是所需变换. $\vec X=0$ 时取恒等.
\end{solution}
\begin{exercise}\label{ps:7V}
\textbf{原题 V.} $n\ge1$, 对两个未来类时向量 $X,Y$, 找共同未来零向量 $L,\overline L$, $m(L,\overline L)=-2$, 及正数 $a,b,c,d$ 使 $X=aL+b\overline L$, $Y=cL+d\overline L$.
\end{exercise}
\begin{solution}
用原题 IV 使 $X=(X^0,0)$, 再空间旋转使 $Y=(Y^0,Y^1,0,\ldots)$; $n=1$ 直接已是此形. 取 $L=(1,1,0,\ldots)$, $\overline L=(1,-1,0,\ldots)$, 则零性及互积 $-2$ 直接成立. $a=b=X^0/2>0$, $c=(Y^0+Y^1)/2>0$, $d=(Y^0-Y^1)/2>0$, 因 $Y^0>|Y^1|$. 逆变换回原系; 所用变换保持未来方向、内积和系数, 即结论.
\end{solution}
\originalsection{作业8: 2011年11月10日交}
\begin{exercise}\label{ps:8I}
\textbf{原题 I.} $T_{\mu\nu}=\phi_\mu\phi_\nu-m_{\mu\nu}Q/2$, $Q=m^{\alpha\beta}\phi_\alpha\phi_\beta$. 若时空梯度非零, 证明对未来类时 $X,Y$ 有 $T(X,Y)>0$.
\end{exercise}
\begin{solution}
补全归一零向量对为 $L,\overline L,e_1,\ldots,e_{n-1}$, 其中横向向量正交单位. 设 $A=L\phi$, $B=\overline L\phi$, $H=\sum(e_i\phi)^2$, 则 $Q=-AB+H$. 因此 $T(L,L)=A^2$, $T(\overline L,\overline L)=B^2$, $T(L,\overline L)=AB+Q=H$. 三者非负, 且若全为0, 则基底所有导数0, 与梯度非零矛盾. 按作业7-V写 $X=aL+b\overline L$, $Y=cL+d\overline L$, 所有系数正,
$T(X,Y)=acA^2+bdB^2+(ad+bc)H>0$. 两者过去指向时同时取负, 双线性使结论不变.
\end{solution}
\begin{exercise}\label{ps:8II}
\textbf{原题 II(a--f), 勘误题.} 三维空间中 $K=(1+t^2+r^2,2tx)$, $J_K^\mu=-T^{\mu\nu}K_\nu$. (a)证未来类时; (b)证 $\partial_\mu K_\nu+\partial_\nu K_\mu=4tm_{\mu\nu}$; (c)原题声称 $\operatorname{tr}_mT=0$; (d)原题声称 $\partial_\mu J_K^\mu=0$; (e)求零方向表达; (f)原题声称 $\int J_K^0$ 守恒. 原卷(d,f)还把波方程印作 $m^{\mu\nu}\phi_\mu\phi_\nu=0$, 此处按题意恢复 $\Box\phi=0$. 其中(c,d,f)不成立, 要说明错误并给修正, 不能证明假命题.
\end{exercise}
\begin{solution}
(a) $K^0>0$, 且 $(K^0)^2-|\vec K|^2=[1+(t+r)^2][1+(t-r)^2]>0$.
(b) $K_0=-(1+t^2+r^2)$, $K_i=2tx_i$. 因此 $2\partial_0K_0=-4t$, $\partial_0K_i+\partial_iK_0=2x_i-2x_i=0$, $\partial_iK_j+\partial_jK_i=4t\delta_{ij}$, 给所需式.
(c) 在4维时空 $m_{\mu\nu}T^{\mu\nu}=Q-4Q/2=-Q$, 一般非零. 例如波解 $\phi=t$ 有 $Q=-1$、迹1.
(d) 场方程给 $\partial_\mu T^{\mu\nu}=0$. 用对称性与(b), 得 $\partial_\mu J_K^\mu=-2t\operatorname{tr}_mT=2tQ$, 一般不为0.
(e) $r>0$ 取 $L=\partial_t+\partial_r$, $\overline L=\partial_t-\partial_r$, $H=|\nabla\phi|^2-\phi_r^2$. 置 $A=1+(t+r)^2$, $B=1+(t-r)^2$, 则 $K=(AL+B\overline L)/2$, $\partial_t=(L+\overline L)/2$. 原题 I 的零标架计算给
\[
J_K^0=\tfrac14\{A(L\phi)^2+B(\overline L\phi)^2+2(1+t^2+r^2)H\}.
\]
$r=0$ 以原笛卡尔表达解释, 不赋予任意径向值.
(f) 紧支撑波解满足 $\frac{\mathrm d}{\mathrm dt}\int J_K^0=2t\int Q$, 不是恒0. 例如非零紧支撑 $f$、$g=0$, 初时 $\int Q=\int|\nabla f|^2>0$, 连续性给小正时刻导数非零. 正确修正流为 $\widetilde J^\mu=J_K^\mu-2t\phi\partial^\mu\phi+\phi^2\partial^\mu t$; 其守恒及正平方表达在下节官方 Bonus 全部证明.
\end{solution}
